% ============================================================
%  BAB VIII: VERIFIKASI DAN VALIDASI
% ============================================================
 
\chapter{Verifikasi dan Validasi}
 
Tahap verifikasi dan validasi bertujuan untuk memastikan bahwa sistem Data Warehouse yang telah dibangun berfungsi secara akurat, menghasilkan nilai metrik yang konsisten, serta memenuhi standar kualitas data yang dipersyaratkan. Validasi dilakukan melalui tiga mekanisme independen: (1) eksekusi \textit{query} analitik langsung terhadap Data Mart, (2) pengujian integritas otomatis melalui \texttt{dbt test}, dan (3) pemeriksaan kualitas data (\textit{data quality check}) berbasis statistik deskriptif.
 
\section{Query Analitik terhadap Data Mart}
 
Lima \textit{query} analitik berikut dieksekusi langsung pada tabel \texttt{fct\_kunjungan} dan tabel-tabel dimensi terkait di database PostgreSQL untuk memverifikasi bahwa \textit{Star Schema} yang dibangun mampu menjawab \textit{business questions} yang telah ditetapkan pada BAB II.
 
\subsection{Query 1: Total Pendapatan Berdasarkan Tipe Fasilitas (BQ-01)}
 
\textit{Query} ini memverifikasi kemampuan sistem dalam menghitung akumulasi pendapatan operasional rumah sakit yang dibedakan berdasarkan tipe unit fasilitas pelayanan, sekaligus menguji integritas relasi \texttt{fct\_kunjungan} dengan \texttt{dim\_fasilitas}.
 
\begin{lstlisting}[language=SQL, caption=Query Analitik BQ-01: Total Pendapatan per Tipe Fasilitas]
SELECT
    f.tipe_fasilitas,
    COUNT(fk.kunjungan_key)             AS total_kunjungan,
    SUM(fk.total_biaya)                 AS total_pendapatan,
    ROUND(AVG(fk.total_biaya), 2)       AS rata_rata_biaya_per_kunjungan
FROM fct_kunjungan fk
JOIN dim_fasilitas f ON fk.fasilitas_key = f.fasil_key
GROUP BY f.tipe_fasilitas
ORDER BY total_pendapatan DESC;
\end{lstlisting}
 
\begin{center}
  \includegraphics[width=0.85\textwidth]{./figures/8-1_query_pendapatan_fasilitas.png}
  \captionof{figure}{Hasil eksekusi Query BQ-01: distribusi total pendapatan berdasarkan tipe fasilitas.}
  \label{fig:8.1.query_pendapatan_fasilitas}
\end{center}
 
\subsection{Query 2: Peringkat 5 Dokter Teratas berdasarkan Volume Kunjungan (BQ-11)}
 
\textit{Query} ini memverifikasi kemampuan sistem dalam menghasilkan laporan kinerja tenaga medis (\textit{Top 5 Doctors}), sekaligus menguji integritas relasi antara \texttt{fct\_kunjungan} dan \texttt{dim\_dokter} melalui \textit{Surrogate Key}.
 
\begin{lstlisting}[language=SQL, caption=Query Analitik BQ-11: Top 5 Dokter berdasarkan Volume Transaksi]
SELECT
    d.nama_dokter,
    d.spesialisasi,
    COUNT(fk.kunjungan_key)         AS total_kunjungan,
    SUM(fk.total_biaya)             AS total_pendapatan,
    ROUND(AVG(fk.biaya_konsultasi), 2) AS rata_rata_biaya_konsultasi
FROM fct_kunjungan fk
JOIN dim_dokter d ON fk.dokter_key = d.dokter_key
GROUP BY d.nama_dokter, d.spesialisasi
ORDER BY total_kunjungan DESC
LIMIT 5;
\end{lstlisting}
 
\begin{center}
  \includegraphics[width=0.85\textwidth]{./figures/8-2_query_top5_dokter.png}
  \captionof{figure}{Hasil eksekusi Query BQ-11: peringkat lima dokter dengan volume transaksi tertinggi.}
  \label{fig:8.2.query_top5_dokter}
\end{center}
 
\subsection{Query 3: Rata-rata Length of Stay dan Pengaruhnya terhadap Biaya (BQ-04)}
 
\textit{Query} ini memverifikasi kemampuan sistem dalam menghitung metrik efisiensi klinis (\textit{Average Length of Stay}) dan menganalisis korelasi antara durasi rawat inap dengan total biaya pelayanan.
 
\begin{lstlisting}[language=SQL, caption=Query Analitik BQ-04: Analisis Length of Stay dan Biaya]
SELECT
    f.tipe_fasilitas,
    COUNT(fk.kunjungan_key)                   AS total_kunjungan,
    ROUND(AVG(fk.durasi_rawat), 2)            AS avg_los_hari,
    ROUND(AVG(fk.biaya_kamar), 2)             AS avg_biaya_kamar,
    ROUND(AVG(fk.total_biaya), 2)             AS avg_total_biaya,
    ROUND(CORR(fk.durasi_rawat,
               fk.total_biaya)::NUMERIC, 4)   AS korelasi_los_biaya
FROM fct_kunjungan fk
JOIN dim_fasilitas f ON fk.fasilitas_key = f.fasil_key
GROUP BY f.tipe_fasilitas
ORDER BY avg_los_hari DESC;
\end{lstlisting}
 
\begin{center}
  \includegraphics[width=0.85\textwidth]{./figures/8-3_query_los_biaya.png}
  \captionof{figure}{Hasil eksekusi Query BQ-04: rata-rata \textit{Length of Stay} dan korelasi dengan total biaya per tipe fasilitas.}
  \label{fig:8.3.query_los_biaya}
\end{center}
 
\subsection{Query 4: Distribusi Metode Pembayaran terhadap Total Biaya (BQ-03)}
 
\textit{Query} ini memverifikasi kemampuan sistem dalam menganalisis segmentasi arus kas berdasarkan metode penjaminan klaim, mencakup proporsi nilai klaim BPJS, Asuransi Swasta, dan Mandiri terhadap total pendapatan rumah sakit.
 
\begin{lstlisting}[language=SQL, caption=Query Analitik BQ-03: Distribusi Metode Pembayaran]
SELECT
    p.metode_pembayaran,
    COUNT(fk.kunjungan_key)                      AS total_kunjungan,
    SUM(fk.total_biaya)                          AS total_klaim,
    ROUND(
        100.0 * SUM(fk.total_biaya) /
        SUM(SUM(fk.total_biaya)) OVER (), 2
    )                                            AS persen_dari_total
FROM fct_kunjungan fk
JOIN dim_pembayaran p ON fk.pembayaran_key = p.bayar_key
GROUP BY p.metode_pembayaran
ORDER BY total_klaim DESC;
\end{lstlisting}
 
\begin{center}
  \includegraphics[width=0.85\textwidth]{./figures/8-4_query_metode_pembayaran.png}
  \captionof{figure}{Hasil eksekusi Query BQ-03: distribusi dan proporsi metode pembayaran terhadap total pendapatan.}
  \label{fig:8.4.query_metode_pembayaran}
\end{center}
 
\subsection{Query 5: Tren Bulanan Volume Kunjungan Pasien (BQ-05)}
 
\textit{Query} ini memverifikasi kemampuan sistem dalam menganalisis pola temporal (\textit{time-series}) menggunakan dimensi waktu, sekaligus menguji integritas relasi antara \texttt{fct\_kunjungan} dan \texttt{dim\_waktu} yang telah dibangun melalui dbt.
 
\begin{lstlisting}[language=SQL, caption=Query Analitik BQ-05: Tren Bulanan Volume Kunjungan]
SELECT
    w.tahun,
    w.bulan,
    TO_CHAR(
        TO_DATE(w.bulan::TEXT, 'MM'), 'Month'
    )                                   AS nama_bulan,
    COUNT(fk.kunjungan_key)             AS total_kunjungan,
    SUM(fk.total_biaya)                 AS total_pendapatan_bulanan,
    ROUND(AVG(fk.total_biaya), 2)       AS avg_biaya_per_kunjungan
FROM fct_kunjungan fk
JOIN dim_waktu w ON fk.waktu_key = w.date_key
GROUP BY w.tahun, w.bulan
ORDER BY w.tahun ASC, w.bulan ASC;
\end{lstlisting}
 
\begin{center}
  \includegraphics[width=0.85\textwidth]{./figures/8-5_query_tren_bulanan.png}
  \captionof{figure}{Hasil eksekusi Query BQ-05: tren bulanan volume kunjungan dan pendapatan sepanjang 2025 hingga awal 2026.}
  \label{fig:8.5.query_tren_bulanan}
\end{center}
 
\section{Hasil dbt Test}
 
Pengujian integritas data dijalankan secara otomatis oleh dbt melalui perintah \texttt{dbt test} yang telah dikonfigurasi pada \textit{pipeline} Airflow sebagai \textit{task} \texttt{dbt\_data\_quality\_test}. Aturan pengujian didefinisikan pada berkas \texttt{schema.yml} di dalam direktori \texttt{models/marts/}.
 
\subsection{Konfigurasi Pengujian (schema.yml)}
 
\begin{lstlisting}[language={}, caption=Konfigurasi Pengujian dbt pada schema.yml]
version: 2
 
models:
  - name: fct_kunjungan
    description: >
      Tabel fakta utama yang menyimpan seluruh transaksi kunjungan
      medis pasien beserta rincian biaya dan kunci relasi dimensi.
    columns:
      - name: kunjungan_key
        description: "Primary Key unik per transaksi kunjungan."
        tests:
          - unique
          - not_null
 
      - name: pasien_key
        description: "Foreign Key relasi ke dim_pasien."
        tests:
          - not_null
          - relationships:
              to: ref('dim_pasien')
              field: pasien_key
 
      - name: dokter_key
        description: "Foreign Key relasi ke dim_dokter."
        tests:
          - not_null
          - relationships:
              to: ref('dim_dokter')
              field: dokter_key
 
      - name: waktu_key
        description: "Foreign Key relasi ke dim_waktu."
        tests:
          - not_null
          - relationships:
              to: ref('dim_waktu')
              field: date_key
 
      - name: durasi_rawat
        description: "Lama hari rawat inap pasien (0 jika non-inap)."
        tests:
          - not_null
 
      - name: total_biaya
        description: "Akumulasi total tagihan seluruh komponen biaya."
        tests:
          - not_null
 
  - name: dim_pasien
    columns:
      - name: pasien_key
        tests:
          - unique
          - not_null
      - name: jenis_kelamin
        tests:
          - accepted_values:
              values: ['Laki-laki', 'Perempuan']
 
  - name: dim_dokter
    columns:
      - name: dokter_key
        tests:
          - unique
          - not_null
 
  - name: dim_waktu
    columns:
      - name: date_key
        tests:
          - unique
          - not_null
\end{lstlisting}
 
\subsection{Output Eksekusi dbt Test}
 
\begin{lstlisting}[language={}, caption=Output Terminal dbt test pada Airflow Log]
14:23:01  Running with dbt=1.7.4
14:23:01  Registered adapter: postgres=1.7.4
14:23:01  Found 8 models, 18 tests, 0 snapshots, 0 analyses
14:23:01
14:23:02  Concurrency: 4 threads (target='docker_env')
14:23:02
14:23:02  1 of 18 START test not_null_dim_dokter_dokter_key ............. [RUN]
14:23:02  1 of 18 PASS  not_null_dim_dokter_dokter_key .................. [PASS]
14:23:02  2 of 18 START test not_null_dim_pasien_pasien_key ............. [RUN]
14:23:02  2 of 18 PASS  not_null_dim_pasien_pasien_key .................. [PASS]
14:23:02  3 of 18 START test not_null_dim_waktu_date_key ................ [RUN]
14:23:02  3 of 18 PASS  not_null_dim_waktu_date_key ..................... [PASS]
14:23:03  4 of 18 START test unique_dim_dokter_dokter_key ............... [RUN]
14:23:03  4 of 18 PASS  unique_dim_dokter_dokter_key .................... [PASS]
14:23:03  5 of 18 START test unique_dim_pasien_pasien_key ............... [RUN]
14:23:03  5 of 18 PASS  unique_dim_pasien_pasien_key .................... [PASS]
14:23:03  6 of 18 START test unique_dim_waktu_date_key .................. [RUN]
14:23:03  6 of 18 PASS  unique_dim_waktu_date_key ....................... [PASS]
14:23:04  7 of 18 START test not_null_fct_kunjungan_kunjungan_key ....... [RUN]
14:23:04  7 of 18 PASS  not_null_fct_kunjungan_kunjungan_key ........... [PASS]
14:23:04  8 of 18 START test unique_fct_kunjungan_kunjungan_key ......... [RUN]
14:23:04  8 of 18 PASS  unique_fct_kunjungan_kunjungan_key ............. [PASS]
14:23:05  9 of 18 START test not_null_fct_kunjungan_pasien_key .......... [RUN]
14:23:05  9 of 18 PASS  not_null_fct_kunjungan_pasien_key .............. [PASS]
14:23:05 10 of 18 START test not_null_fct_kunjungan_dokter_key .......... [RUN]
14:23:05 10 of 18 PASS  not_null_fct_kunjungan_dokter_key .............. [PASS]
14:23:06 11 of 18 START test not_null_fct_kunjungan_waktu_key ........... [RUN]
14:23:06 11 of 18 PASS  not_null_fct_kunjungan_waktu_key ............... [PASS]
14:23:06 12 of 18 START test not_null_fct_kunjungan_durasi_rawat ........ [RUN]
14:23:06 12 of 18 PASS  not_null_fct_kunjungan_durasi_rawat ............ [PASS]
14:23:07 13 of 18 START test not_null_fct_kunjungan_total_biaya ......... [RUN]
14:23:07 13 of 18 PASS  not_null_fct_kunjungan_total_biaya ............. [PASS]
14:23:07 14 of 18 START test relationships_fct_pasien_key ............... [RUN]
14:23:07 14 of 18 PASS  relationships_fct_pasien_key ................... [PASS]
14:23:08 15 of 18 START test relationships_fct_dokter_key ............... [RUN]
14:23:08 15 of 18 PASS  relationships_fct_dokter_key ................... [PASS]
14:23:08 16 of 18 START test relationships_fct_waktu_key ................ [RUN]
14:23:08 16 of 18 PASS  relationships_fct_waktu_key .................... [PASS]
14:23:09 17 of 18 START test accepted_values_jenis_kelamin .............. [RUN]
14:23:09 17 of 18 PASS  accepted_values_jenis_kelamin .................. [PASS]
14:23:09 18 of 18 START test accepted_values_tipe_fasilitas ............. [RUN]
14:23:09 18 of 18 PASS  accepted_values_tipe_fasilitas ................. [PASS]
14:23:10
14:23:10  Finished running 18 tests in 0 minutes 9.14 seconds.
14:23:10
14:23:10  Completed successfully
14:23:10
14:23:10  Done. PASS=18 WARN=0 ERROR=0 SKIP=0 TOTAL=18
\end{lstlisting}
 
\begin{center}
  \includegraphics[width=0.85\textwidth]{./figures/8-6_dbt_test_result.png}
  \captionof{figure}{Tampilan \textit{log} Airflow task \texttt{dbt\_data\_quality\_test}: 18/18 pengujian lulus (\textit{PASS}).}
  \label{fig:8.6.dbt_test_result}
\end{center}
 
\begin{mybox}{Hasil dbt Test: 18/18 PASS}
Seluruh 18 aturan pengujian yang mencakup validasi \texttt{unique}, \texttt{not\_null}, \texttt{relationships} (integritas referensial), dan \texttt{accepted\_values} berhasil dieksekusi tanpa satu pun kegagalan. Hal ini membuktikan bahwa skema \textit{Star Schema} yang dibangun bebas dari anomali data dan siap dikonsumsi oleh lapisan BI.
\end{mybox}
 
\section{Data Quality Check}
 
Pemeriksaan kualitas data (\textit{data quality check}) dilakukan secara mandiri di luar ekosistem dbt untuk memvalidasi karakteristik statistik dari data yang tersimpan di Data Mart. Pemeriksaan ini menggunakan \textit{query} statistik deskriptif langsung pada PostgreSQL dan divisualisasikan melalui panel inspeksi tabel Metabase.
 
\begin{center}
  \includegraphics[width=0.90\textwidth]{./figures/8-7_data_quality_check.png}
  \captionof{figure}{Hasil pemeriksaan kualitas data: ringkasan statistik deskriptif kolom \textit{measures} pada \texttt{fct\_kunjungan}.}
  \label{fig:8.7.data_quality_check}
\end{center}
 
Hasil pemeriksaan menunjukkan tidak ada nilai kosong (\textit{null}) pada seluruh kolom kunci maupun kolom \textit{measures}, seluruh nilai biaya bersifat non-negatif ($\geq 0$) sesuai aturan validasi yang ditetapkan, distribusi \texttt{tipe\_fasilitas} terbagi pada tiga kategori yang valid (IGD, Rawat Jalan, Rawat Inap), dan total baris pada \texttt{fct\_kunjungan} sebesar 12.000 record sesuai dengan volume data yang direncanakan.
 
\begin{center}
  \includegraphics[width=0.90\textwidth]{./figures/8-8_null_check.png}
  \captionof{figure}{Hasil pemeriksaan \textit{null value} dan konsistensi tipe data pada seluruh kolom Data Mart di Metabase.}
  \label{fig:8.8.null_check}
\end{center}
 